\section{Discussion and Limitations}
\label{sec:discussion}

Future prediction and reusable representation require separate evidence. Clinical-status, progression, report-fidelity, and probability-reliability scores assess forecasting, while the three downstream tasks assess information accessible through the current state. The generic output slots describe available observations without guaranteeing complete, interpretable, or disentangled clinical states. Downstream adaptation does not by itself establish that longitudinal training improves these representations. All three task heads can exploit the direct image-token pathway. Matched image-only controls are therefore necessary to establish any contribution from $S_t$; task accuracy alone does not show that the state retains more useful information.

Visual slots remain defined when the report is withheld; performance with and without clinical text requires evaluation. Predicting $\widehat S_{t,h}$ does not establish image synthesis, full-report generation, or future dense outputs from decoders that also require the corresponding input image. Medical JEPA retraining and transfer of volumetric knowledge to radiographs remain untested directions.

The model learns observational follow-ups. Follow-up selection, treatment, care setting, scanner changes, and historical report content can correlate with recorded futures. Clinical-record timestamp filtering and target-only future evidence constrain the forecasting inputs without identifying natural history or causal treatment effects. Current reports are assumed available retrospectively because reliable signing times are unavailable. Conclusions depend on the evaluated cohorts, horizons, and annotations. Supplementary retrieval additionally depends on candidate-pool size and eligibility; their coverage must be reported.

VLM adaptation makes training cost consequential. Stop-gradient isolates target outputs from the prediction loss but does not ensure stability or prevent collapse. The EMA target and simultaneous current-task supervision are specified, but their effects on training stability and representation quality still require empirical validation. Future reports must remain unavailable to the prediction inputs.

Zero-shot comparisons combine task adaptation, output-format learning, and slot conditioning. The matched no-slots arm isolates conditioning within the shared decoder, but retains temporal and visual consistency objectives. Native Qwen fine-tuning, removal of temporal supervision, and alternative pooled-feature conditions are needed to distinguish these effects. The 300-question VQA subset, strict JSON scoring, and single-seed estimates limit the current evidence; full-test evaluation and uncertainty estimates remain pending.
